\subsection{Common collision bookkeeping}
\label{app:lower-proof-collisions}

We count how matching edges leave the unresolved-edge set; this pathwise
identity supplies the batch-size constraint in the collision lower-tail
proof. Recall that \(\UnresolvedVertices{t}\) contains the vertices on
edges still unresolved immediately before commit \(t\), after any preceding
queries. At this boundary, operation \(s_t\), the unresolved edge set is
\(\OpUnresolvedMatching{s_t}\), as in
\cref{app:eq:unresolved-matching-update,app:eq:remaining-trigger-set}.
Put \(\UnresolvedVertexCount{t}:=|\UnresolvedVertices{t}|\).
Removing whole edges makes this number even.

For the selected batch \(B_t\), define its active part and collision count:
\begin{equation}
 \ActiveBatch{t}:=B_t\cap\UnresolvedVertices{t},
 \qquad
 \ActiveBatchSize{t}:=|\ActiveBatch{t}|,
 \qquad
 \CollisionCount{t}:=
 \bigl|\{e\in\OpUnresolvedMatching{s_t}:
               e\subseteq\ActiveBatch{t}\}\bigr|,
 \qquad
 \TotalCollisions:=\sum_t\CollisionCount{t}.
 \label{app:eq:unresolved-batch-and-collisions}
\end{equation}
An unresolved edge has no committed endpoint, so
\(\UnresolvedVertices{t}\subseteq\ResidualPositionsAt{G_{t-1}}\).
The inclusion can be strict: query-retired endpoints and the uncommitted
mate of a commit-retired endpoint remain uncommitted but are no longer
active. Figure~\ref{fig:c4-matching-states}(b) shows the latter case.
A later commit of that mate creates no new unresolved-edge collision.

Let \(\ActiveRoundCount:=|\{t:\ActiveBatchSize{t}\ge1\}|\le\RoundBudget{}\)
be the number of rounds with a nonempty active batch.

\begin{lemma}[Pathwise edge-counting identity]
\label{app:lem:pathwise-edge-counting}
Under the retirement rules in \cref{app:eq:unresolved-matching-update},
every complete execution satisfies the pathwise identity
\begin{equation}
 \frac N2
 =\GenieDeletions+
  \sum_t(\ActiveBatchSize{t}-\CollisionCount{t}).
 \label{app:eq:lower-edge-ledger}
\end{equation}

Consequently,
\begin{equation}
 \sum_{t:\ActiveBatchSize{t}\ge1}
   (\ActiveBatchSize{t}-1)
 =\frac N2-\GenieDeletions+\TotalCollisions-\ActiveRoundCount
 \ge\frac N2-\GenieDeletions-\RoundBudget{}.
 \label{app:eq:active-excess-ledger}
\end{equation}
\end{lemma}

\begin{proof}
An active batch touches exactly
\(\ActiveBatchSize{t}-\CollisionCount{t}\) previously unresolved edges.
Indeed, its \(\CollisionCount{t}\) internal edges use
\(2\CollisionCount{t}\) vertices, while every remaining active vertex lies
on a distinct crossing edge. Thus the number touched is
\[
 \underbrace{\CollisionCount{t}}_{\text{internal edges}}
 +\underbrace{(\ActiveBatchSize{t}-2\CollisionCount{t})}_{\text{crossing edges}}
 =\ActiveBatchSize{t}-\CollisionCount{t}.
\]
Since every position is eventually committed, every original matching
edge is retired exactly once: either at a successful counterfactual
trigger test (query retirement), or at its first active commit batch
(commit retirement). Summing over the \(N/2\) original edges proves
\eqref{app:eq:lower-edge-ledger}.

Finally,
\(\sum_{t:\ActiveBatchSize{t}\ge1}(\ActiveBatchSize{t}-1)
=\sum_t\ActiveBatchSize{t}-\ActiveRoundCount\).
Substituting
\(\sum_t\ActiveBatchSize{t}=N/2-\GenieDeletions+\TotalCollisions\)
from \eqref{app:eq:lower-edge-ledger}, then using
\(\TotalCollisions\ge0\) and \(\ActiveRoundCount\le\RoundBudget{}\),
gives \eqref{app:eq:active-excess-ledger}.
\end{proof}
